% GENERATED FILE. Edit canonical lessons, then run npm run generate:books.
% canonical-source-hash: fnv1a64:121841bd8bb56db0
% canonical-lessons: AR-C291-idha, AR-C291-idha-amtarat, AR-R291-if-it-rains

\chapter{If It Rains}
\label{ch:ar-if-it-rains}
\hlchaptermodality{291}

\begin{chapteropening}
\textbf{By the end of this chapter:} \emph{I can make a plan that depends on a real condition, and give its result.}

Complete AR-R291-if-it-rains: a real condition with idhā and the past, and its result with fa-.
\end{chapteropening}

\section[\texorpdfstring{\ar{إذا} (idhā)}{إذا (idhā)}]{\ar{إذا} (idhā) — if}

\label{lesson:AR-C291-idha}



\begin{quote}
Say \textquotedblleft{}I have been working here for three years\textquotedblright{}, then \textquotedblleft{}I arrived a week ago\textquotedblright{}. (\emph{Aʿmalu hunā mundhu thalāth sanawāt}, \emph{waṣaltu qabla usbūʿ}.)

You have known \emph{idhā}, \textquotedblleft{}if\textquotedblright{}, since chapter two hundred and seven. What it takes after it will surprise you.
\end{quote}

\begin{grammarlens}[title={Grammar Lens: idhā and the past}]
After \textbf{\ar{إذا}} — \emph{idhā} — Arabic puts a \textbf{past} verb, and means the future:

\begin{quote}
\textbf{\ar{إذا} \ar{وصلت} \ar{مبكرا،} \ar{أشرب} \ar{قهوة}.} — \emph{idhā waṣaltu mubakkiran, ashrabu qahwa} — If I arrive early, I'll have a coffee.
\end{quote}

\emph{Waṣaltu} looks like \textquotedblleft{}I arrived\textquotedblright{}, but after \emph{idhā} it is \textquotedblleft{}if I arrive\textquotedblright{}. The condition is something that may well happen: you are not sure, but you expect it. The second half, the result, can be a plain present.

With \emph{kāna}: \emph{idhā kāna ʿindī waqt, adhhabu maʿaka}, if I have time, I'll go with you.
\end{grammarlens}

\begin{grammarlens}[title={Grammar Lens: in, the other \textquotedblleft{}if\textquotedblright{}}]
\textbf{\ar{إن}} — \emph{in} — is \textquotedblleft{}if\textquotedblright{} as well, usually described as a little less sure than \emph{idhā}. It takes the past in the same way: \emph{in waṣalta mubakkiran}, if you arrive early. You hear it inside a phrase you may know already: \emph{in shāʾa allāh}, \textquotedblleft{}if God wills\textquotedblright{}.

\emph{In}, if, with no doubled letter, is a different word from \emph{anna}, that, and \emph{an}, to.
\end{grammarlens}

\subsection*{Your turn}
Say these aloud:

\begin{itemize}
\raggedright
  \item \textquotedblleft{}If I arrive early, I'll have a coffee\textquotedblright{} — \emph{idhā waṣaltu mubakkiran, ashrabu qahwa}
  \item \textquotedblleft{}If I have time, I'll go with you\textquotedblright{} — \emph{idhā kāna ʿindī waqt, adhhabu maʿaka}
  \item \textquotedblleft{}if you arrive early\textquotedblright{}, with \emph{in} — \emph{in waṣalta mubakkiran}
  \item a plan of your own with \emph{idhā} and a past verb, then a present
\end{itemize}

\begin{tcolorbox}[breakable,colback=teal!4,colframe=teal!35!black,title={Before you move on}]
After \emph{idhā}, which tense? (\textbf{The past}, meaning the future.) \textquotedblleft{}If I arrive\textquotedblright{}? (\emph{\textbf{Idhā waṣaltu.}}) The other \textquotedblleft{}if\textquotedblright{}? (\emph{\textbf{In}}, as in \emph{in shāʾa allāh}.)
\end{tcolorbox}
\section[\texorpdfstring{\ar{إذا} \ar{أمطرت} (idhā amṭarat)}{إذا أمطرت (idhā amṭarat)}]{\ar{إذا} \ar{أمطرت} (idhā amṭarat) — if it rains}

\label{lesson:AR-C291-idha-amtarat}



\begin{quote}
Say \textquotedblleft{}if I arrive early, I'll have a coffee\textquotedblright{}. (\emph{Idhā waṣaltu mubakkiran, ashrabu qahwa}.) Then \textquotedblleft{}if you arrive early\textquotedblright{}, with the other \textquotedblleft{}if\textquotedblright{}. (\emph{In waṣalta mubakkiran}.)

The result in that sentence was a plain present. Many results are not, and then Arabic marks the seam.
\end{quote}

\begin{grammarlens}[title={Grammar Lens: fa- before the result}]
\begin{quote}
\textbf{\ar{إذا} \ar{أمطرت،} \ar{فسنشاهد} \ar{فيلما} \ar{في} \ar{المنزل}.} — \emph{idhā amṭarat, fa-sa-nushāhidu fīlman fī al-manzil} — If it rains, we'll watch a film at home.
\end{quote}

\emph{Amṭarat} is \textquotedblleft{}it rained\textquotedblright{}, and after \emph{idhā}, \textquotedblleft{}it rains\textquotedblright{}. The result opens with \emph{sa-}, and so it takes \textbf{\ar{فـ}} in front: \emph{fa-sa-nushāhidu}. It is the \emph{fa-} of \textquotedblleft{}so\textquotedblright{}: the condition comes true, so this follows.
\end{grammarlens}

\begin{grammarlens}[title={Grammar Lens: four results that take fa-}]
A result needs \emph{fa-} when it opens with \emph{sa-} or \emph{sawfa}, with \emph{lan}, with an order, or with a noun. Four sentences, one of each:

\begin{quote}
\emph{idhā amṭarat, fa-sa-nushāhidu fīlman} · \emph{idhā amṭarat, fa-lan nakhruja}, we won't go out · \emph{idhā waṣalta, fa-ttaṣil bī}, call me · \emph{idhā taʾakhkhara al-qiṭār, fa-l-ḥāfila afḍal}, the bus is better
\end{quote}

A plain present, like \emph{ashrabu qahwa}, needs none. When you are unsure, the \emph{fa-} before \emph{sa-} is the one you will use most.
\end{grammarlens}

\subsection*{Your turn}
Say these aloud:

\begin{itemize}
\raggedright
  \item \textquotedblleft{}If it rains, we'll watch a film at home\textquotedblright{} — \emph{idhā amṭarat, fa-sa-nushāhidu fīlman fī al-manzil}
  \item \textquotedblleft{}If it rains, we won't go out\textquotedblright{} — \emph{idhā amṭarat, fa-lan nakhruja}
  \item \textquotedblleft{}If you arrive, call me\textquotedblright{} — \emph{idhā waṣalta, fa-ttaṣil bī}
  \item \textquotedblleft{}If the train is late, the bus is better\textquotedblright{} — \emph{idhā taʾakhkhara al-qiṭār, fa-l-ḥāfila afḍal}
  \item your own plan for tomorrow: \emph{idhā} … \emph{fa-sa-} …
\end{itemize}

\begin{tcolorbox}[breakable,colback=teal!4,colframe=teal!35!black,title={Before you move on}]
\textquotedblleft{}If it rains\textquotedblright{}? (\emph{\textbf{Idhā amṭarat.}}) What goes in front of a result with \emph{sa-}? (\emph{\textbf{Fa-.}}) Name the four openings that call for it. (\emph{\textbf{Sa-}} \textbf{or \emph{sawfa}, \emph{lan}, an order, a noun.}) Does a plain present need it? (\textbf{No.})
\end{tcolorbox}
\section[(review)]{Review — if it rains}

\label{lesson:AR-R291-if-it-rains}



\begin{quote}
Say \textquotedblleft{}the train is better than the bus\textquotedblright{} and \textquotedblleft{}the biggest city\textquotedblright{}. (\emph{Al-qiṭār afḍal min al-ḥāfila}, \emph{akbar madīna}.) Then \textquotedblleft{}for two years\textquotedblright{} and \textquotedblleft{}a week ago\textquotedblright{}. (\emph{Mundhu sanatayn}, \emph{qabla usbūʿ}.)
\end{quote}

\begin{grammarlens}[title={What you've built}]
Two lessons, and a condition with both its halves. The condition: \emph{idhā} or \emph{in} with a past verb, meaning something that may well happen. The result: a plain present as it is, or \emph{fa-} in front when it opens with \emph{sa-}, \emph{lan}, an order or a noun.

A plan with its condition in one breath:

\begin{quote}
\textbf{\ar{إذا} \ar{كان} \ar{الطقس} \ar{جميلا،} \ar{فسنذهب} \ar{إلى} \ar{البحر}.} — \emph{idhā kāna aṭ-ṭaqs jamīlan, fa-sa-nadhhabu ilā al-baḥr} — If the weather is nice, we'll go to the sea.
\end{quote}
\end{grammarlens}

\subsection*{Your turn}
Say these aloud:

\begin{itemize}
\raggedright
  \item \emph{idhā kāna aṭ-ṭaqs jamīlan, fa-sa-nadhhabu ilā al-baḥr}
  \item \textquotedblleft{}If you get lost, phone me\textquotedblright{} — \emph{idhā ḍalalta, fa-ttaṣil bī}
  \item \textquotedblleft{}If the pharmacist phones, I'll give him the number\textquotedblright{} — \emph{idhā ittaṣala aṣ-ṣaydalī, fa-sa-uʿṭīhi ar-raqm}
  \item \textquotedblleft{}If I learn a lot, I'll become an engineer\textquotedblright{} — \emph{idhā taʿallamtu kathīran, fa-sa-uṣbiḥu muhandisan}
  \item \textquotedblleft{}If it rains, I won't travel\textquotedblright{} — \emph{idhā amṭarat, fa-lan usāfira}
  \item two plans of your own for the weekend: one with \emph{idhā} and a plain present, one with \emph{fa-sa-}
\end{itemize}

\begin{tcolorbox}[breakable,colback=teal!4,colframe=teal!35!black,title={Before you move on}]
\textquotedblleft{}If it rains, we won't go out\textquotedblright{}? (\emph{\textbf{Idhā amṭarat, fa-lan nakhruja.}}) Which tense follows \emph{idhā}? (\textbf{The past.}) Which \textquotedblleft{}if\textquotedblright{} sits in \emph{in shāʾa allāh}? (\emph{\textbf{In.}})
\end{tcolorbox}
